\subsection*{The remaining quartic CM cohomology groups at five}
\begin{statement}{E.112}{Proved}{Certified character fields and complete central five-primary fibres}
Put $R=\mathbb Z_5$, $S=\{2,5,\infty\}$,
$V=T_5E_+$ for $E_+:y^2=x^3+x$, and retain the actual
rank-two direct summand $U=U_{\rm CM}\subset\operatorname{Sym}^2V$
and its projection and section from (E108.24).
Let $\chi_2$ be the quadratic character of $\mathbb Q(\sqrt2)$.
For each of the following five lattices,
\[
 \mathcal N=\{V,\ V\chi_2,\ V(1),\ V(1)\chi_2,\ U\},
 \tag{E112.1}
\]
the finite and continuous groups over the specified arithmetic source are
\[
 \begin{gathered}
 H^0(G_{\mathbb Q,S},N/5^n)=H^2(G_{\mathbb Q,S},N/5^n)=0,\qquad
 H^1(G_{\mathbb Q,S},N/5^n)\simeq\mathbb Z/5^n,\quad n\ge1,\\
 H^1_{\rm cts}(G_{\mathbb Q,S},N)\simeq R,\qquad
 H^2_{\rm cts}(G_{\mathbb Q,S},N)=0\qquad(N\in\mathcal N).
 \end{gathered}
 \tag{E112.2}
\]
The proof constructs compatible generators $h_N$ and their inverse
scalar-coordinate maps.  Over the unrestricted arithmetic source,
\[
 \begin{gathered}
 H^1_{\rm cts}(G_{\mathbb Q},V)
 \simeq H^1_{\rm cts}(G_{\mathbb Q},V\chi_2)
 \simeq H^1_{\rm cts}(G_{\mathbb Q},V(1))
 \simeq H^1_{\rm cts}(G_{\mathbb Q},V(1)\chi_2)\simeq R,\\
 H^1_{\rm cts}(G_{\mathbb Q},U)
       \simeq R\oplus\mathcal R_{\mathcal P_-}(R),
 \qquad \mathcal P_-=\{p\text{ prime}:p\equiv3\pmod4\},\\
 \mathcal R_P(R)=\{(t_p)\in\prod_{p\in P}R:
       \#\{p:t_p\notin5^nR\}<\infty\text{ for every }n\}.
 \end{gathered}
 \tag{E112.3}
\]
All isomorphisms include their actual continuous maps in both
directions, as specified below.  Thus for the affine plus-quartic
module $A=W\oplus V^3$ of E.103, the full central module of every
nonempty intrinsic class-two fibre is
\[
 \begin{split}
 H^1_{\rm cts}(G_{\mathbb Q},\Lambda^2 A)\simeq{}&R^{66}
   \oplus\mathcal R_{\mathcal P_{\mathbb Q}}(R)^6\\
 &\oplus\mathcal R_{\{p:p\equiv1\ (4)\}}(R)^3
   \oplus\mathcal R_{\mathcal P_-}(R)^3 .
 \end{split}
 \tag{E112.4}
\]
Here $\mathcal P_{\mathbb Q}$ is the set of all rational primes.
The first two restricted factors have the valuation coordinates
of E.108.  The last factors have the literal tame-inertia
coordinates (E112.17) below.  The $66$ scalar coordinates have
the ordered decomposition $30+15+18+3$ inherited from (E108.7).
The maps to and from the original central coefficient module
are exactly the character columns and CM maps of E.108.
The central action on sections and the inverse central ratio are
those of E.105 and E.108.
For $R=\mathbb Q_5$, replace every free $R$ by $\mathbb Q_5$
and each restricted module by the space of $\mathbb Q_5$
sequences tending to zero.  These formulas evaluate the fibre
groups. E.115 evaluates the mixed obstruction deciding the image
of the full quartic abelian section space with these coefficients.
\end{statement}

\begin{proof}
\emph{Literal division fields and both field inclusions.}
Let $K=\mathbb Q(i)$, and for $j=1,2$ put
\[
 \begin{gathered}
 \sigma_1=-1,\quad\sigma_2=1,\qquad
 a_j^4=\sigma_j(1+2i),\quad b_j^4=\sigma_j(1-2i),\qquad
 L_j=K(a_j,b_j),\\
 E_1:y^2=x^3+x,\qquad E_2:y^2=x^3+4x.
 \end{gathered}
 \tag{E112.5}
\]
Choose the radicals so that complex conjugation sends $a_j$ to
$b_j$.  The first radical polynomial is Eisenstein at the
Gaussian prime $1+2i$.  Its discriminant is supported at $2$
and that prime, so the other Gaussian prime $1-2i$ is unramified
in $K(a_j)$.  At any prime of $K(a_j)$ above $1-2i$ the element
$\sigma_j(1-2i)$ has valuation one.  The second radical polynomial
is therefore Eisenstein there.  This proves
$[L_j:K]=16$ and $[L_j:\mathbb Q]=32$.
The following automorphisms are all distinct and give the entire
relative Galois group:
\[
 \begin{gathered}
 g_{r,s}(a_j)=i^ra_j,\quad g_{r,s}(b_j)=i^sb_j,
       \quad(r,s)\in(\mathbb Z/4)^2,\\
 c(i)=-i,\quad c(a_j)=b_j,\quad c(b_j)=a_j,
       \qquad c g_{r,s}c^{-1}=g_{-s,-r}.
 \end{gathered}
 \tag{E112.6}
\]
They also prove that $L_j/\mathbb Q$ is Galois.  These fields are
unramified outside $2$ and $5$, directly from the radical
polynomial discriminants and $\operatorname{disc}(K)=-4$.

Here are actual torsion points.  For $E_1$ define
\[
 \begin{aligned}
 y_+&=\frac{1+i}{2-i}a_1,&
 x_+&=\frac{y_+^2}{(4+2i)/5},\\
 y_-&=\frac{1-i}{2+i}b_1,&
 x_-&=\frac{y_-^2}{(4-2i)/5}.
 \end{aligned}
 \tag{E112.7}
\]
For $E_2$ use instead
\[
 \begin{aligned}
 y_+&=\frac4{2-i}a_2,&
 x_+&=\frac{y_+^2}{4(4+2i)/5},\\
 y_-&=\frac4{2+i}b_2,&
 x_-&=\frac{y_-^2}{4(4-2i)/5}.
 \end{aligned}
 \tag{E112.8}
\]
In each case $P_\pm=(x_\pm,y_\pm)$ satisfies the indicated
curve equation and
\[
 [2]P_+=[i]P_+,\qquad [2]P_-=[-i]P_-,\qquad
             [5]P_+=[5]P_-=0,
 \quad [i](x,y)=(-x,iy).
 \tag{E112.9}
\]
These identities follow by substituting (E112.5) into the
duplication formulas
$x([2]P)=((3x^2+A)/(2y))^2-2x$ and
$y([2]P)=((3x^2+A)/(2y))(x-x([2]P))-y$ for $A=1,4$.
The actual exact field calculation verifies every substitution.
For example in (E112.7) the identity $x_+=a_1^{-2}$ gives
$x_+^2=(-1+2i)/5$; the duplication ordinate is $2x_+^3/y_+$,
which equals $iy_+$.  The same substitution with
$x_+^2=4(-1+2i)/5$ proves the $A=4$ identity.
Squaring the first two endomorphism identities gives
$[4]P_\pm=[-1]P_\pm$, hence the displayed fifth-torsion relation.
Both ordinates are nonzero.  The two $\mathbb F_5$-eigenvalues
of $[i]$ on these points are $2$ and $-2$, so the points are
independent and form a basis of $E_j[5]$.

Consequently $\mathbb Q(E_j[5])\subset L_j$.  Conversely its
coordinates contain
$i=y([2]P_+)/y(P_+)$, and (E112.7) or (E112.8) recovers both
$a_j$ and $b_j$ by multiplication by the displayed nonzero
elements of $K$.  Thus
\[
 L_j=\mathbb Q(E_j[5]),\qquad
 \mu_5\subset L_j,
 \qquad
 T_5E_1\otimes\chi_2\simeq T_5E_2.
 \tag{E112.10}
\]
The middle inclusion uses the primitive Weil pairing of the
two basis points.  The last map is induced by the actual
isomorphism $(x,y)\mapsto(2x,2\sqrt2\,y)$; its inverse is
$(X,Y)\mapsto(X/2,Y/(2\sqrt2))$.  The sign under
$\sqrt2\mapsto-\sqrt2$ proves precisely the asserted twist.

Fix the root $i_5\in\mathbb Z_5$ of $i_5^2=-1$ with
$i_5\equiv2\pmod5$.  The integral CM projectors
$(1\pm I/i_5)/2$ split $V$ into two rank-one lattices over
$G_K$.  Choose bases reducing to $P_+,P_-$.
On $g_{r,s}$ the two residual Tate characters are
$\alpha=2^r$ and $\beta=2^{-s}$, because $[i]P_-=[-2]P_-$.
In particular the cyclotomic character is $2^{r-s}$ and
the two characters of $U(-1)/5$ are
\[
             \alpha/\beta=2^{r+s},\qquad
             \beta/\alpha=2^{-r-s}.
 \tag{E112.11}
\]
Their common kernel fixes $K(a_1b_1)=K(5^{1/4})$.
The distinct character $2^{r-s}$ fixes the subfield
corresponding to the vector $(1,3)$ in the next calculation.
Thus the passage from the actual CM contraction in E.108 to
the field used below keeps its full character, including the
sign of the second exponent.

\emph{Six certified subfields and explicit ideal descent.}
For each $j$ and each vector
\[
 \mathcal V=\{(1,0),(0,1),(1,1),(1,2),(1,3),(2,1)\},
 \quad C_{u,v}=\{(r,s):ur+vs=0\pmod4\},
 \tag{E112.12}
\]
the field $F_{j,u,v}=K(a_j^ub_j^v)$ is the fixed field
$L_j^{C_{u,v}}$, of degree eight over $\mathbb Q$ and
degree four over $K$.  This follows both from the character
in (E112.6) and from the literal minimal polynomials certified
in the calculation.  The six kernels exhaust the kernels of
the twelve primitive characters, paired under multiplication
by $-1$.  Every one of the sixteen characters of
$(\mathbb Z/4)^2$ is trivial on at least one $C_{u,v}$.

The following is the complete certified class-number table.
The second and third numeric columns are absolute field
discriminants; all class numbers in the last column are the
same in the two families.
\[
\begin{array}{c|r|r|c}
(u,v)&|d(F_{1,u,v})|&|d(F_{2,u,v})|&h(F_{j,u,v})\\\hline
(1,0)&32768000&8192000&1\\
(0,1)&32768000&8192000&1\\
(1,1)&64000000&64000000&2\\
(1,2)&819200000&204800000&2\\
(1,3)&4000000&4000000&1\\
(2,1)&819200000&204800000&2
\end{array}
\tag{E112.13}
\]
Here is the reproducible certificate, including the field
maps used to apply this table.  The program constructs a
compositum of $x^8-2\sigma_jx^4+5$ and $x^4-5$.
PARI's returned coordinate polynomials give $a_j,t_j$ in
its degree-$32$ quotient, with $t_j^4=5$; set
$b_j=t_j/a_j$ and $i=(a_j^4/\sigma_j-1)/2$.
It checks the primitive-generator identity, (E112.5), and
all point identities above in that exact quotient.
For $m=a_j^ub_j^v$, write $m^4=A+Bi$ with integers $A,B$.
For $B\ne0$ the degree-eight polynomial is
$x^8-2Ax^4+A^2+B^2$, with the literal map $x\mapsto m$
and inverse $i\mapsto(x^4-A)/B$ in the subfield.
For $(u,v)=(1,1)$ the program uses the exact compositum
of $x^4-5$ and $x^2+1$, retains both coordinate polynomials,
and checks their images in $L_j$; its polynomial is
$x^8+4x^6-4x^4+64x^2+16$.
Every minimal polynomial is checked irreducible, and every
displayed field inclusion and copy of $K$ is checked by
substitution before the class-group calculation.
Each complete $\operatorname{bnfinit}(P,1)$ is certified by
$\operatorname{bnfcertify}$ with its default flag $0$, with
return value $1$ in all twelve cases.  The complete structures
and all embedding polynomials are retained in the binary
certificate; the numeric records include the class invariants
and decompositions of $2$ and $5$.  The certification meaning
is the original PARI specification
\cite[\texttt{bnfcertify}, default flag $0$]{E_PARICertifiedUnits}.

To transfer this finite result, put $A_j=\operatorname{Cl}(L_j)\{5\}$.
For every indicated $C=C_{u,v}$ the actual ideal extension
and norm maps satisfy
\[
 \begin{gathered}
 j_C:\operatorname{Cl}(L_j^C)\{5\}\longrightarrow A_j,
 \quad [\mathfrak a]\longmapsto[\mathfrak a\mathcal O_{L_j}],\\
 N_C:A_j\longrightarrow\operatorname{Cl}(L_j^C)\{5\},
 \quad[\mathfrak b]\longmapsto[N_{L_j/L_j^C}\mathfrak b],\\
 N_Cj_C=[4],\qquad j_CN_C=\sum_{c\in C}c,
 \qquad j_C:\operatorname{Cl}(L_j^C)\{5\}\xrightarrow{\sim}A_j^C,
 \quad j_C^{-1}=\tfrac14N_C.
 \end{gathered}
 \tag{E112.14}
\]
The two compositions follow on prime ideals from their
ramification and residue degrees, and therefore on all ideals
and their classes.  They prove injectivity and surjectivity
onto $A_j^C$ because multiplication by four is invertible
on a finite five-group.  Thus (E112.13) gives $A_j^C=0$
for all six kernels.  More explicitly, the sixteen integral
idempotents
$e_\psi=16^{-1}\sum_g\psi(g)^{-1}g\in R[(\mathbb Z/4)^2]$
are pairwise orthogonal and sum to one.  For every $\psi$
choose a listed $C$ on which $\psi$ is trivial; then
$e_\psi A_j\subset A_j^C=0$.  This proves
\[
            \operatorname{Cl}(L_j)\{5\}=0,
       \qquad\operatorname{Pic}(\mathcal O_{L_j}[1/10])/5=0.
 \tag{E112.15}
\]
The second group is a quotient of the ordinary ideal class
group.  No value for the full class number of either
degree-$32$ field is needed or asserted.

\emph{The local decomposition maps kill the finite obstruction group.}
For the three subfields used by the residual characters of
$N(-1)$, the complete prime decompositions have the following
ramification and residue degrees.  They are identical for
$j=1,2$:
\[
\begin{array}{c|c|c|c|c}
(u,v)&(e,f)\text{ over }2&(e,f)\text{ over }5&
 |D_{2}(F_{j,u,v}/K)|&|D_{5}(F_{j,u,v}/K)|\\\hline
(1,0),(0,1)&(8,1)&(4,1),(1,4)&4&4,4\\
(1,1)&(4,2)&(4,1),(4,1)&4&4,4
\end{array}
\tag{E112.16}
\]
The relative order at $2$ is the absolute local degree
divided by $[K_2:\mathbb Q_2]=2$.  At $5$ both primes of
$K$ have completion $\mathbb Q_5$.  Thus every indicated
decomposition group maps onto the cyclic quotient of order
four.  Every faithful character of that quotient has zero
fixed subspace.  This proves the local assertion for both
characters of $V/5$ using $(1,0),(0,1)$, for their inverses
using the same two kernels, and for both characters of
$U(-1)/5$ using $(1,1)$ and (E112.11).

For $N=V,V(1),U$ use $L=L_1$; for their $\chi_2$ twists
use $L=L_2$ and (E112.10).  In every case $L$ trivializes
$N/5$ and $\mu_5$.  Restriction identifies
$H^2(G_{\mathbb Q,S},N/5)$ with
$H^2(G_{L,S_L},N/5)^{\Delta}$, where
$\Delta=\operatorname{Gal}(L/\mathbb Q)$ has order $32$.
Indeed positive-degree cohomology of $\Delta$ on every
five-primary module is killed by the averaging homotopy;
restriction and $32^{-1}$ times corestriction are inverse
on these invariants.  This records both directions of descent.

The Galois--etale comparison identifies the group over $L$
with cohomology of $\operatorname{Spec}\mathcal O_L[1/10]$.
The finite Kummer sequence and (E112.15) give
$H^2(\mathcal O_L[1/10],\mu_5)=
\operatorname{Br}(\mathcal O_L[1/10])[5]$.
Evaluation of homomorphisms supplies the actual coefficient
map
$\mu_5\otimes (N/5)(-1)\to N/5$ and its inverse; a chosen
primitive root and a basis give mutually inverse coordinate
matrices.  Applied to cochains, it identifies the group with
$\operatorname{Br}(\mathcal O_L[1/10])[5]\otimes (N/5)(-1)$.
The Brauer invariant map injects its first factor into
$\bigoplus_{w\mid2,5}\mathbb F_5$, with its natural
permutation action.  The field has no real places.
After tensoring and taking exact $\Delta$-invariants, evaluation
at one prime over each rational prime identifies the invariant
permutation terms with
$((N/5)(-1))^{D_2}\oplus((N/5)(-1))^{D_5}$.
The inverse sends a fixed vector to its translates at the
other primes; independence of the coset representative is
exactly the fixed-vector condition.  Both terms are zero
by (E112.16).  This proves finite degree-two vanishing modulo
five.  The comparison, Kummer, and Brauer maps used here are
the exact maps in \cite[II.2.1, II.2.2(a), II.2.9 and proof]{E_MilneADT}.

The residual $G_K$ characters of $V,V(1),U$, in exponent
coordinates $(r,s)$, are respectively
$\{(1,0),(0,-1)\}$, $\{(2,-1),(1,-2)\}$, and
$\{(2,0),(0,-2)\}$.  All are nontrivial; the twist case
uses its actual $E_2$ field with the same list.
Thus $H^0(G_{\mathbb Q,S},N/5)=0$.
The short exact coefficient sequences
$0\to N/5^{n-1}\xrightarrow{5}N/5^n\to N/5\to0$
give vanishing of $H^0$ and $H^2$ for every $n$ by induction.
All these finite cohomology groups are finite.
Complex conjugation on each lattice in (E112.1) has
one positive and one negative integral eigenline: for $U$
it exchanges the two pure CM squares, and for the Tate
twists multiplication by $-1$ exchanges the two signs.
The finite global Euler formula therefore gives
$|H^1(G_{\mathbb Q,S},N/5^n)|=5^n$
\cite[I.4.15; I.5.1 and proof]{E_MilneADT}.
The reduction maps on $H^1$ are onto by the degree-two
vanishing for $N/5$.  The image of multiplication by five
on the group at level $n$ is the included group at level
$n-1$: factor multiplication through reduction to that
level and use its surjectivity.  Thus the quotient by five
has order five.  The group of order $5^n$ is consequently
cyclic, as asserted in (E112.2).

Choose a nonzero class at level one and successively lift
it along these surjections.  Each lift generates its cyclic
group.  Choose representing cocycles; at each step their
reductions differ by a coboundary, and lifting its zero-cochain
to the next coefficient module corrects that difference.
This constructs a compatible system of actual cocycles and
its continuous limit $h_N:G_{\mathbb Q,S}\to N$.
The map $a\mapsto[a h_N]$ has inverse that reads, at every
level $n$, the unique coefficient in $\mathbb Z/5^n$ of
the chosen generator.  The inverse coefficients are compatible.
These are continuous inverse maps for the inverse-limit
topology.  Finiteness of the level-one cohomology groups at
every exponent kills the derived inverse limit in degree two;
finiteness of degree-zero groups does the same in degree one
\cite[Proposition 1.6 and Theorem 2.2]{E_JannsenContinuous}.
This proves both continuous assertions in (E112.2).

\emph{Every unrestricted cocycle and its prime coordinates.}
At a prime $p\notin S$ the coefficient action is unramified.
Restriction of an integral cocycle to inertia is a homomorphism
from its tame $\mathbb Z_5(1)$ quotient, and its value belongs
to $N(-1)^{F_p=1}$, where $F_p$ is arithmetic Frobenius.
Wild inertia contributes zero because it is pro-$p$.
For $N=V(1),V(1)\chi_2$ the eigenvalues on $N(-1)$ have
complex absolute value $\sqrt p$; for $N=V,V\chi_2$ they
have absolute value $1/\sqrt p$.  None equals one, so the
inertia map is zero.  A cocycle with zero restrictions to
all these inertia groups descends through their closed
normal closure to $G_{\mathbb Q,S}$.  Inflation and this
descent are inverse, including coboundaries.  This proves
the first line of (E112.3).

For $N=U$, choose the CM bases $e_+,e_-$ just fixed.
At $p\equiv1\pmod4$, Frobenius preserves the two lines,
with eigenvalues $\alpha_p,\beta_p$ and
$\alpha_p\beta_p=p$.  On $U(-1)$ the eigenvalues are
$\alpha_p/\beta_p,\beta_p/\alpha_p$.
Equality to one would give equal roots of
$T^2-a_pT+p$, hence $a_p^2=4p$ with integer $a_p$,
which is impossible for a prime.  Thus the integral fixed
module at every such prime is zero.
At $p\equiv3\pmod4$, the relation $F_pI=-IF_p$ gives
an actual matrix $\left(\begin{smallmatrix}0&a_p'\\b_p'&0\end{smallmatrix}\right)$
on $V$, with $a_p',b_p'\in R^\times$ and
$a_p'b_p'=-p$ by the Weil determinant.
Use the induced basis of $U(-1)$, including its chosen
Tate basis.  Its Frobenius, fixed-line section, and inverse
coordinate are
\[
 \begin{gathered}
 F_p^U=\begin{pmatrix}0&(a_p')^2/p\\(b_p')^2/p&0\end{pmatrix},
 \quad w_p=\begin{pmatrix}1\\(b_p')^2/p\end{pmatrix},
 \quad\ell_p=\tfrac12\bigl(1,(a_p')^2/p\bigr),\\
 (F_p^U)^2=1,\qquad F_p^Uw_p=w_p,\qquad
 \ell_pw_p=1,\qquad w_p\ell_p=(1+F_p^U)/2,\\
 t_p(z)=\ell_p\bigl(z|_{I_p}(\tau_p)\bigr),
 \qquad z|_{I_p}(\tau_p)=t_p(z)w_p.
 \end{gathered}
 \tag{E112.17}
\]
Here the value on a compatible tame generator is interpreted
in $U(-1)$, so a different choice of Tate basis changes
both displayed maps by its actual unit.  The entries
$a_p',b_p'$ are defined by Frobenius on the finite groups
$E_+[5^n]$, compatibly in $n$, and are therefore specified
by the original elliptic coefficient representation.
The four matrix identities follow from $a_p'b_p'=-p$.
They prove exactly that the fixed module is $Rw_p$ with
inverse row $\ell_p$.

For each $p\in\mathcal P_-$, localization on
$\operatorname{Spec}\mathbb Z[1/10]$ at $p$ gives the
finite-coefficient exact maps
$H^1(G_{\mathbb Q,S},U/5^n)\to
H^1(G_{\mathbb Q,S\cup\{p\}},U/5^n)\to
H^0(\mathbb F_p,(U/5^n)(-1))\to
H^2(G_{\mathbb Q,S},U/5^n)$.
The last term is zero.  Purity identifies the middle map
with the actual tame-inertia value just used
\cite[II.0.1(f), II.1.7(b), II.2.9]{E_MilneADT}.
The fibres over $w_p\bmod5^n$ are nonempty finite sets.
Their inverse limit is nonempty: all finite compatibility
conditions can be met by choosing a lift at their largest
exponent, and compactness of the product of the finite
sets supplies a compatible choice.
The preceding zero-cochain correction gives an actual
integral cocycle $\xi_p$ unramified outside $S\cup\{p\}$
whose inertia value is $w_p$.  Fix one such cocycle for
each prime.  Its values lie in the integral lattice, so
all these cocycles have the common coefficient bound one.

The required unrestricted forward map is now explicit:
\[
 \begin{gathered}
 R\oplus\mathcal R_{\mathcal P_-}(R)
       \longrightarrow H^1_{\rm cts}(G_{\mathbb Q},U),\\
 (a,(t_p))\longmapsto
       \left[a h_U+\sum_{p\in\mathcal P_-}t_p\xi_p\right].
 \end{gathered}
 \tag{E112.18}
\]
The sum converges uniformly as a sum of integral continuous
cochains, because $t_p\to0$ and the chosen cochains are
uniformly bounded.  At $q$ its inertia value is precisely
$t_qw_q$.
Conversely the reduction of any continuous cocycle modulo
$5^n$, together with the finite coefficient action, defines
a finite extension of $\mathbb Q$.  That extension is
ramified at only finitely many primes.  Consequently its
coordinates $t_p$ from (E112.17) belong to
$\mathcal R_{\mathcal P_-}(R)$.
Subtract the actual convergent sum $\sum t_p\xi_p$.
The resulting cocycle has zero inertia outside $S$ and
therefore descends to $G_{\mathbb Q,S}$.
Read its unique scalar coefficient $a$ by (E112.2).
This is the inverse of (E112.18) in both directions.
An integral coboundary restricts to zero at the primes in
question, so the inverse is defined on cohomology classes.
The residue rows and the convergent-sum map are bounded
for the coefficient congruence topology; the final scalar
map is the finite inverse-limit coordinate already proved
continuous.  This proves the topological assertion as well.

Finally substitute these maps into the actual decomposition
(E108.7).  Its peripheral part contributes $30$ free
coordinates, its mixed parts contribute $15+18$, and its
three copies of $U$ contribute three more.  Its six
determinant and three twisted determinant factors retain
their original valuation maps.  The three new residue
factors are exactly (E112.17)--(E112.18).  These explicit
coefficient sections and projections sum to the identity
on $\Lambda^2A$ by E.108, proving (E112.4) and both fibre
coordinate maps with no additional kernel or quotient.

A continuous cochain on the compact domain
$G_{\mathbb Q}^r$ with values in $N\otimes\mathbb Q_5$
has bounded image, and thus lies in $5^{-m}N$ for some
single $m$.  The rational cochain complex is therefore
the union of these scaled integral complexes.
Filtered unions are exact, so its cohomology is the
integral cohomology tensored with $\mathbb Q_5$.
A sequence tending to zero in $\mathbb Q_5$ has a common
denominator bound, giving exactly the rational modules
specified after (E112.4).

The program checks every field inclusion and torsion
identity, all twelve full class certificates, all sixteen
character-kernel covers, the full residual Galois group
and its conjugation law, the Fourier and norm compositions
at six coefficient exponents, all local decomposition
orders, and the polynomial identities of (E112.17).
Its finite matrix examples supplement the polynomial
proof; the finite-S cohomology, inverse-limit, and
unrestricted support arguments are the preceding proof.
The earlier interrupted degree-$32$ certification and
the two rejected GP parser attempts are retained with
their actual outputs and statuses.  None enters (E112.13).
\end{proof}
